\section{Methodology}
\label{submission:method}

In this section, we formally outline the model merging problem, describe the unsupervised test-time optimization paradigm of AdaMerging \cite{AdaMerging_ICLR_2024} across task-wise and layer-wise formulations, define the mathematical constructs of our two diagnostic treatments, and formalize the theoretical framework for the **Overfitting-Optimizer Paradox** and the newly discovered **Spatial Averaging Paradox**.

\subsection{Preliminaries and Model Merging Setup}
\label{sub:preliminaries}
Let $\theta_0 \in \mathbb{R}^D$ denote the parameter weight vector of a pre-trained base neural network (e.g., a Vision Transformer visual encoder). Suppose we fine-tune this base model independently on $T$ distinct classification tasks, yielding $T$ specialized expert models, denoted by $\theta_1, \dots, \theta_T \in \mathbb{R}^D$. 

For each task $t \in \{1, \dots, T\}$, we compute its corresponding \emph{task vector} (or parameter update) $\tau_t$ by subtracting the pre-trained weights from the fine-tuned weights:
\begin{equation}
\tau_t = \theta_t - \theta_0
\end{equation}
The objective of model merging is to integrate these task-specific capabilities into a single multi-task network $\theta_{\text{merged}}$ without retraining on any task's original training data. The classic baseline, Task Arithmetic \cite{task_arithmetic}, performs a simple linear combination of these task vectors:
\begin{equation}
\theta_{\text{merged}}(\lambda_{\text{static}}) = \theta_0 + \sum_{t=1}^T \lambda_{\text{static}} \tau_t
\end{equation}
where $\lambda_{\text{static}} > 0$ is a globally applied static scaling factor (typically chosen via grid search on validation splits).

Following standard practice in multi-task visual model merging \cite{task_arithmetic, AdaMerging_ICLR_2024}, the model merging process is performed solely on the shared visual encoder weights ($\theta_0$ and $\theta_t$), while the specialized classification heads are kept task-specific, disjoint, and unmerged. During evaluation, input images are routed to their respective task-specific heads to compute predictions. This setup implicitly assumes test-time knowledge of the task identity (referred to as \emph{Oracle Routing}) to select the correct classification head. This evaluation protocol is crucial for isolating and measuring the functional quality of the merged representation encoder independently of classification head conflicts.

\subsection{Unsupervised Test-Time AdaMerging}
Rather than using static scales, AdaMerging \cite{AdaMerging_ICLR_2024} parameterizes the merging process. In the most expressive, \textbf{layer-wise formulation}, each layer $l \in \{1, \dots, L\}$ of each task $t \in \{1, \dots, T\}$ is assigned a unique scaling coefficient $\Lambda_{l, t} \in \mathbb{R}$. The parameters of the merged model at layer $l$ are given by:
\begin{equation}
\theta_{\text{merged}}^{(l)}(\Lambda) = \theta_0^{(l)} + \sum_{t=1}^T \Lambda_{l, t} \tau_{l, t}
\end{equation}
In the more restricted, \textbf{task-wise formulation}, exactly $T$ coefficients $\lambda = [\lambda_1, \dots, \lambda_T]^T \in \mathbb{R}^T$ are optimized, assigning a single global scale per task across all layers:
\begin{equation}
\theta_{\text{merged}}^{(l)}(\lambda) = \theta_0^{(l)} + \sum_{t=1}^T \lambda_t \tau_{l, t}
\end{equation}
To learn these coefficients without training data, AdaMerging leverages unsupervised test-time adaptation. Given a small unlabeled calibration batch $X$ drawn from the joint test-time multi-task distribution, the coefficients ($\Lambda$ or $\lambda$) are optimized by minimizing the Shannon entropy of the model's predictions. Crucially, prediction entropy minimization is fundamentally a transductive, unsupervised surrogate objective. Unlike supervised cross-entropy, prediction entropy does not measure functional alignment or class-conditional correctness; instead, it purely measures the confidence (sharpness) of prediction distributions on the specific batch $X$. Consequently, it is a highly noisy and unconstrained objective that is susceptible to pathological shortcuts—such as scaling up weights to artificially inflate logit magnitudes—especially when optimized under global weight-space bottlenecks without task labels. Under the joint test-time prediction entropy objective, the loss is defined as:
\begin{equation}
\label{eq:entropy_loss}
\mathcal{L}(\Theta; X) = -\frac{1}{|X|} \sum_{x \in X} \sum_{c=1}^C p_c(x, \Theta) \log p_c(x, \Theta)
\end{equation}
where $p_c(x, \Theta) = p_c\left(x; \theta_{\text{merged}}(\Theta)\right)$ is the predicted softmax probability of class $c \in \{1, \dots, C\}$ for input sample $x$ under the merged model, and $\Theta \in \{\Lambda, \lambda\}$.

\subsection{Diagnostic Treatments of the Overfitting-Optimizer Paradox}
Because layer-wise AdaMerging optimizes hundreds of variables ($\Lambda \in \mathbb{R}^{L \times T}$) on tiny calibration datasets, we hypothesize that it is susceptible to test-time overfitting. This is the \textbf{Overfitting-Optimizer Paradox}: while the optimizer successfully captures a functional, high-dimensional representational routing (where layer-wise parameters are structurally specialized to their positions in the neural network hierarchy), they also capture high-frequency, transductive test-time overfitting dynamics from the prediction entropy objective. To deconstruct these scaling dynamics, we define two diagnostic control treatments applied to optimized layer-wise coefficients:

\textbf{1. Intra-Task Layer Shuffling.}
For each task $t$, let $\{\Lambda_{l, t}\}_{l=1}^L$ be the set of optimized layer-wise coefficients. We define a random permutation function $\pi_t: \{1, \dots, L\} \rightarrow \{1, \dots, L\}$. We then reconstruct the merged model using these permuted coefficients:
\begin{equation}
\label{eq:shuffling}
\theta_{\text{shuffled}}^{(l)} = \theta_0^{(l)} + \sum_{t=1}^T \Lambda_{\pi_t(l), t} \tau_{l, t}
\end{equation}
By shuffling the layer-specific scale coefficients, we break the correspondence between each learned coefficient and its specific position in the network hierarchy. If the learned parameters are highly specialized to their respective representational layers, shuffling will severely disrupt this structural alignment, causing a collapse in performance. This allows us to diagnose whether the coefficients are uniformly generalizable or deeply tied to the network's architectural hierarchy.

\textbf{2. Spatial Averaging (Spatial Mean).}
We take the optimized layer-wise coefficients and compute their flat spatial average (mean) over all layers for each task:
\begin{equation}
\bar{\lambda}_t = \frac{1}{L} \sum_{l=1}^L \Lambda_{l, t}
\end{equation}
We then construct a regularized merged model using this single flat scalar per task across all layers:
\begin{equation}
\label{eq:averaging}
\theta_{\text{averaged}}^{(l)} = \theta_0^{(l)} + \sum_{t=1}^T \bar{\lambda}_t \tau_{l, t}
\end{equation}
This post-hoc spatial averaging acts as a powerful regularizer, reducing the parameter degrees of freedom back to a flat task-level scale. By doing so, it serves as a spatial low-pass filter that smooths away the high-frequency test-time overfitting component of individual layers, while preserving the robust task-level global scales.

\subsection{The Spatial Averaging Paradox and Multi-Task Gradient Imbalance}
Our deconstruction reveals a secondary, more profound optimization anomaly, which we term the \textbf{Spatial Averaging Paradox}:
\begin{equation}
\text{Acc}\left(\theta_{\text{averaged}}\right) \gg \text{Acc}\left(\theta_{\text{merged}}(\lambda^*)\right)
\end{equation}
where $\lambda^*$ is the direct task-wise optimized scaling vector. That is, \emph{indirect} optimization (optimizing $L \times T$ parameters and averaging them post-hoc) succeeds, whereas \emph{direct} optimization (optimizing exactly $T$ global parameters) fails spectacularly, even degrading performance below its uniform initialization.

We formalize and explain this paradox through **multi-task gradient imbalance** on uncalibrated test-time prediction entropy. Note that the Shannon prediction entropy in Equation~\ref{eq:entropy_loss} is extremely sensitive to logit scaling. For any task $t$, scaling up its task vector $\tau_t$ by a factor $\lambda_t > 1$ directly inflates the logit magnitudes of the network. This artificially sharpens the softmax distribution, driving the prediction entropy of task $t$ to near-zero:
\begin{equation}
\lim_{\lambda_t \to \infty} p\left(x; \theta_0 + \lambda_t \tau_t\right) \in \{0, 1\}^C \implies \lim_{\lambda_t \to \infty} \mathcal{L}_{\text{task } t} = 0
\end{equation}
Importantly, prediction entropy is not calibrated across tasks of varying difficulty. Simple classification tasks (such as MNIST) have sharp, high-confidence logit distributions, whereas harder tasks (such as CIFAR-10 or SVHN) have naturally flatter, low-confidence distributions. 

Under direct task-wise optimization in a shared, low-dimensional bottleneck $\lambda \in \mathbb{R}^T$, the gradient of the joint entropy loss is dominated by easy tasks because their entropy is highly responsive to scaling:
\begin{equation}
\left\| \nabla_{\lambda_{\text{easy}}} \mathcal{L}_{\text{easy}} \right\| \gg \left\| \nabla_{\lambda_{\text{hard}}} \mathcal{L}_{\text{hard}} \right\|
\end{equation}
The global optimizer is thus drawn into a pathologic region where it massively scales up easy-task coefficients to drive joint entropy down. Under a low-dimensional bottleneck, this joint scaling creates severe parameter interference in shared weights (such as early self-attention projections), collapsing performance on the harder tasks.

Conversely, under high-dimensional layer-wise optimization ($\Lambda \in \mathbb{R}^{L \times T}$), the optimizer is granted $L \times T$ degrees of freedom. This allows it to adjust layer-wise scales locally (e.g., modifying late layers of easy tasks while preserving early layers of hard tasks) to minimize joint entropy without resorting to global scaling trade-offs. 

When we take the spatial mean of these high-dimensional coefficients post-hoc (Equation~\ref{eq:averaging}), we act as a spatial low-pass filter. This process smooths away the individual layer-wise transductive overfitting component while preserving the regularized global task-wise scaling signal, explaining why post-hoc Spatial Averaging successfully recovers and generalizes.

\subsection{Calibrated Prediction Entropy: A Proposed Algorithmic Remedy}
\label{sub:calibrated_remedy}
Since the root cause of the Spatial Averaging Paradox is the uncalibrated nature of joint prediction entropy and the resulting gradient imbalance across tasks, we propose and evaluate a simple yet elegant algorithmic remedy: \textbf{Calibrated Prediction Entropy}.

Let $\mathcal{L}_t(\Theta; X)$ denote the average prediction entropy on task $t \in \{1, \dots, T\}$. Let $\Theta^{(0)}$ represent the baseline uniform task-scale coefficient vector initialized to $0.3$ (equivalent to standard Task Arithmetic). We define the calibrated test-time adaptation loss as:
\begin{equation}
\mathcal{L}_{\text{calibrated}}(\Theta; X) = \sum_{t=1}^T \frac{\mathcal{L}_t\left(\Theta; X\right)}{\mathcal{L}_t\left(\Theta^{(0)}; X\right) + \epsilon}
\end{equation}
where $\epsilon > 0$ is a small constant (e.g., $10^{-6}$) to ensure numerical stability.

By dividing each task's prediction entropy by its initial entropy at the uniform initialization, we normalize each task's contribution to have an initial loss value of exactly $1.0$. This prevents easy tasks with sharp logit distributions (such as MNIST) from having disproportionately responsive gradients and dominating the shared, low-dimensional optimization landscape. In Section~\ref{submission:experiments}, we empirically evaluate whether this calibration remedy can overcome the Spatial Averaging Paradox and restore direct flat-coefficient adaptation.
